% !TEX root = ../main.tex
% Introduction
% Claims to support: headline claim + RQs + contributions

\section{Introduction}
\label{sec:01_intro}
Issue Tracking Systems provide a structured way to communicate concerns related to the software project by allowing users to submit issue reports~\cite{colavito2024leveraging}. Issue reports can be classified with labels to help with triage. Effective issue report classification (IRC) results in faster issue resolution~\cite{liao2018exploring}, and is crucial for many software maintenance tasks, such as issue prioritization~\cite{catolino2019not}, issue assignment~\cite{wu2022spatial, jeong2009improving}, and developer onboarding~\cite{santos2023tag, steinmacher2018let, aracena2025applying}. However, studies show that labels are rarely applied in issue creation~\cite{antoniol2008bug, herzig2013s, cabot2015exploring, kallis2019ticket}. As a result, project maintainers must manually classify the issues for triage, which is laborious and time-consuming~\cite{fan2017road}.

To address this problem, numerous automated methods have been proposed around supervised learning over three established issue report categories: bug, feature, and question~\cite{antoniol2008bug, kallis2019ticket, izadi2022catiss, aracena2025applying, aracena2024applyinglargelanguagemodels}. Early work relied on classic machine learning (ML) techniques~\cite{antoniol2008bug, kallis2019ticket, fan2017road, herzig2013s}, while more recent studies fine-tune transformer models on labeled issue reports~\cite{izadi2022catiss, izadi2022predicting, trautsch2022predicting, bharadwaj2022github, colavito2022issue, wang2021well, gomes2023bert}. These approaches require a substantial amount of training data and often struggle with minority classes~\cite{izadi2022catiss, izadi2022predicting, colavito2022issue, heo2025study, vargovich2023givemelabeledissues}. The current state-of-the-art methods fine-tune Large Language Models (LLMs) using Low-Rank Adaptation (LoRA~\cite{hu2022lora})~\cite{heo2025study, aracena2025applying, aracena2024applyinglargelanguagemodels}. However, fine-tuning LLMs is computationally expensive and must be repeated whenever the model is swapped or to incorporate new labeled issues that accumulate in the project~\cite{fan2017road, kallis2019ticket}.

Prior work shows that placing a few demonstration examples in the prompt can boost an LLM's performance on software-related tasks~\cite{assi2026llm, le2023log} through in-context learning~\cite{brown2020language}. Since few-shot prompting enhances performance without model training or weight updates, it has the potential to serve as a cost-efficient alternative to fine-tuning LLMs for IRC. Few-shot performance, however, depends on the choice of in-context examples. Studies show that using the query's retrieved nearest neighbors as in-context examples yields strong few-shot performance~\cite{liu2022makes}. Additionally, the nearest neighbors tend to share the same label as the query~\cite{cover1967nearest}, making them suitable in-context examples. The combination of these insights leads to retrieval-augmented few-shot prompting, where the classified issue reports most similar to a query are retrieved and used as its in-context examples. Yet both a systematic evaluation of this approach for IRC and its comparison against fine-tuning remain underexplored. We aim to address both gaps in this study.

An important factor in retrieval-augmented few-shot performance is the number of provided examples: too few may not provide enough context and too many can add noise~\cite{liu2022makes}. Since the in-context examples are the query's nearest neighbors, they can classify the query under the standard $k$-nearest-neighbor (KNN) framework~\cite{cover1967nearest, dudani1976distance}. Therefore, the point where adding neighbors stops improving the retrieval-only classification informs an upper bound for the number of in-context examples. Additionally, retrieval-only IRC provides a non-parametric baseline that isolates the LLM's contribution from that of retrieval in retrieval-augmented few-shot prompting. To our knowledge, no prior study has explored such a baseline for IRC. Therefore, we pose our first research question:

\smallskip
\noindent\textbf{RQ1.\ How effectively can retrieval alone classify issue reports, and at what point does adding more retrieved neighbors stop helping?} We answer this question by applying a standard distance-weighted KNN classification~\cite{cover1967nearest, dudani1976distance}. Given a query issue, its top-$k$ most similar issue reports (nearest neighbors) are first retrieved from an embedded index of existing issue reports. Next, the label is predicted with a similarity-weighted vote over the neighbors. We evaluate this method over a wide range of top-$k$ values to measure its peak performance and identify the point where additional neighbors are not beneficial (best $k$ value). With the best $k$ value established, we ask:

\smallskip
\noindent\textbf{RQ2.\ How well does retrieval-augmented few-shot prompting classify issue reports, and how does performance vary with the number of in-context examples?} We answer this question by presenting the same top-$k$ nearest issue reports as classified examples in a few-shot prompt, instructing the model to predict a label for the query issue. $k{=}0$ results in a zero-shot LLM baseline similar to prior work~\cite{colavito2025benchmarking}. We evaluate this approach with different model scales over a range of $k$ values bounded by RQ1.


RQ2 evaluations reveal that retrieval-augmented few-shot has the weakest performance on the question class, and question issues are frequently misclassified as bugs. Zero-shot prompting results indicate that the models are already biased towards question-to-bug misclassifications even before the examples are presented. We hypothesize that bug-labeled neighbors retrieved for true-question queries reinforce this misclassification. This motivates our next research question:

\smallskip
\noindent\textbf{RQ3.\ Can a balanced retrieval intervention reduce question-to-bug misclassification and improve retrieval-augmented few-shot performance?} We answer this question by exploring a more balanced retrieval technique that excludes bug-labeled neighbors when they outnumber question-labeled ones by a margin in the top-$k$. We evaluate it over the same $k$ values and model scales used for RQ2.

For brevity, we refer to the retrieval-only classifier (RQ1) as \textbf{Vo}ting-based \textbf{TAG}ging (\votag), the retrieval-augmented few-shot approach (RQ2) as \textbf{R}etrieval-\textbf{A}ugmented \textbf{G}enerative \textbf{TAG}ging (\ragtag), and its balanced variant (RQ3) as \textbf{B}alanced \ragtag (\bragtag). We compare these methods against the established LoRA~\cite{hu2022lora} fine-tuning baseline~\cite{heo2025study, aracena2025applying, aracena2024applyinglargelanguagemodels} to answer our final research question:


\smallskip
\noindent\textbf{RQ4.\ How do retrieval-based methods compare against LoRA fine-tuning for IRC in classification performance and computational cost?} We conduct this comparison across the Qwen2.5-Instruct model family at 3B, 7B, 14B, and 32B parameters~\cite{yang2024qwen25} to analyze the effect of model scale without architecture variance. We use Heo et al.'s~\cite{heo2025study} 11-project dataset for our experiments and adopt its two evaluation settings: a \emph{project-specific} (PS) configuration where retrieval/training is over each project's individual data split in isolation, and a \emph{project-agnostic} (PA) configuration where retrieval/training is over data from the 11 projects pooled together.


Results show that retrieved few-shot examples consistently improve IRC performance. Even a single retrieved example outperforms the zero-shot baseline across all model sizes (\Cref{sec:ragtag}). LLM reasoning also consistently improves over pure retrieval: every \ragtag\ configuration outperforms \votag's best result. Notably, however, \votag\ comes within 0.018 macro $F_1$ of Qwen-3B's zero-shot (\Cref{sec:rq1}).

Additionally, retrieval-augmented few-shot prompting is competitive with fine-tuning. Using only each project's own labeled issues (PS), \ragtag\ outperforms fine-tune by 0.021--0.040 macro $F_1$ on all model sizes (\Cref{sec:finetune-eval}). Fine-tune only catches up when trained on $11\times$ more labeled data via cross-project pooling (PA), and even then \ragtag\ remains competitive: aggregated across all model sizes, fine-tune PA leads \ragtag\ PS by 0.029 macro $F_1$, and the two methods are statistically tied at Qwen-3B and Qwen-32B. \bragtag\ closes this gap by reducing the question-to-bug misclassification rate without sacrificing performance on other labels (\Cref{sec:bragtag}). The aggregate ($\bragtag - $ fine-tune) macro $F_1$ difference is $-0.010$. Applying a \votag\ fallback to invalid LLM outputs across all methods tightens this to statistical equivalence, passing two one-sided tests (TOST) at $\delta{=}0.01$.

Retrieval-augmented few-shot offers deployment advantages as well. Averaged across the four model sizes, \ragtag\ and \bragtag\ require 33\% less peak GPU memory than fine-tuning (\Cref{sec:finetune-eval}). Additionally, incorporating newly labeled issues and swapping the underlying model mid-production do not require retraining cycles.

Overall, this study demonstrates retrieval-augmented few-shot's practicality in IRC, particularly when hardware resources are constrained, labeled data is limited, or the deployment requires frequent model or data updates.

The main contributions of this paper are:

\begin{itemize}
  \item Establishing a retrieval-only, non-parametric baseline for IRC that also serves as a fallback mechanism for invalid LLM predictions (\votag).

  \item Conducting the first systematic evaluation of retrieval-augmented
  few-shot prompting for IRC across different LLM scales and numbers of
  in-context examples (\ragtag).

  \item Proposing a balanced retrieval technique for few-shot prompting that
  reduces question-to-bug misclassifications without degrading performance on
  the other labels (\bragtag).

  \item Empirically comparing these proposed methods against the
  established LoRA fine-tuning baseline on both classification performance and
  computational cost, in two data scope settings (PS and PA). We find that
  retrieval-augmented few-shot prompting is competitive with fine-tuning while
  using 11$\times$ less labeled data per project and 33\% less peak GPU memory.
\end{itemize}

The remainder of this paper is organized as follows: \Cref{sec:02_related} reviews prior work. \Cref{sec:03_approach} explains our proposed methods and the fine-tuning baseline. \Cref{sec:setup} details the experimental setup. \Cref{sec:evaluations} explains our
evaluations and reports the results. \Cref{sec:discussion} discusses failure modes, practical implications, and future work. \Cref{sec:threats} addresses the potential threats to validity. Finally, \Cref{sec:conclusion} concludes the paper.


